% =========================================================================
% Capítulo 0: Repaso
% Tema 01: Álgebra de vectores
% Problema 04: Base Esférica
% =========================================================================
\subsubsection{Base Esférica}
Se define la base esférica $\{\vec{e}_\rho, \vec{e}_\phi, \vec{e}_\theta\}$ respecto a la cartesiana:
$$\vec{e}_\rho = \sin\phi\cos\theta\vec{e}_x + \sin\phi\sin\theta\vec{e}_y + \cos\phi\vec{e}_z$$
$$\vec{e}_\phi = \cos\phi\cos\theta\vec{e}_x + \cos\phi\sin\theta\vec{e}_y - \sin\phi\vec{e}_z$$
$$\vec{e}_\theta = -\sin\theta\vec{e}_x + \cos\theta\vec{e}_y$$

\begin{multicols}{2}
\begin{enumerate}
    \item $\vec{e}_\rho \cdot \vec{e}_z =$
    \item $\vec{e}_\phi \cdot \vec{e}_z =$
    \item $\vec{e}_\theta \cdot \vec{e}_z =$
    \item $\vec{e}_\rho \cdot \vec{e}_\rho =$
    \item $\vec{e}_\rho \times \vec{e}_\phi =$
    \item $\vec{e}_\theta \times \vec{e}_\rho =$
    \item Exprese $\vec{e}_z$ en base esférica.
    \item $\vec{e}_\rho \cdot (\sin\phi\cos\theta\vec{e}_x + \sin\phi\sin\theta\vec{e}_y) =$
    \item $\| \vec{e}_\rho \times \vec{e}_z \| =$
    \item $(\vec{e}_\rho \cdot \vec{e}_x)^2 + (\vec{e}_\rho \cdot \vec{e}_y)^2 + (\vec{e}_\rho \cdot \vec{e}_z)^2 =$
\end{enumerate}
\end{multicols}
% --- SOLUCIÓN ---------------------------------------------------------------
\begin{solucion}{Base Esférica}
\begin{enumerate}
    \item $\vec{e}_\rho \cdot \vec{e}_z = (\sin\phi\cos\theta\vec{e}_x + \sin\phi\sin\theta\vec{e}_y + \cos\phi\vec{e}_z) \cdot \vec{e}_z = \cos\phi$.
    \item $\vec{e}_\phi \cdot \vec{e}_z = (\cos\phi\cos\theta\vec{e}_x + \cos\phi\sin\theta\vec{e}_y - \sin\phi\vec{e}_z) \cdot \vec{e}_z = -\sin\phi$.
    \item $\vec{e}_\theta \cdot \vec{e}_z = 0$ (el vector $\vec{e}_\theta$ está contenido en el plano $xy$).
    \item $\vec{e}_\rho \cdot \vec{e}_\rho = 1$ (vector unitario).
    \item $\vec{e}_\rho \times \vec{e}_\phi = \vec{e}_\theta$ (base orto-normal derecha).
    \item $\vec{e}_\theta \times \vec{e}_\rho = \vec{e}_\phi$ (orden cíclico de la base).
    \item $\vec{e}_z = (\vec{e}_z \cdot \vec{e}_\rho)\vec{e}_\rho + (\vec{e}_z \cdot \vec{e}_\phi)\vec{e}_\phi = \cos\phi\vec{e}_\rho - \sin\phi\vec{e}_\phi$.
    \item $\vec{e}_\rho \cdot (\sin\phi\cos\theta\vec{e}_x + \sin\phi\sin\theta\vec{e}_y) = \sin^2\phi\cos^2\theta + \sin^2\phi\sin^2\theta = \sin^2\phi$.
    \item $\| \vec{e}_\rho \times \vec{e}_z \| = \sin(\phi) \cdot 1 = \sin\phi$.
    \item $(\vec{e}_\rho \cdot \vec{e}_x)^2 + (\vec{e}_\rho \cdot \vec{e}_y)^2 + (\vec{e}_\rho \cdot \vec{e}_z)^2 = \sin^2\phi\cos^2\theta + \sin^2\phi\sin^2\theta + \cos^2\phi = 1$.
\end{enumerate}
\end{solucion}
